% sections/A_data.tex: data collection, curation, reproducibility and notation (appendix).
% Owner: the appendix-A writer. Static content handed over from the skeleton: fig:gesture (docs/img frames),
% fig:system (TikZ, updated to the Stage 2 data flow), tab:sensors (browser sensors) and tab:notation.
% Every result number is a generated key (\R...). The thresholds typed here (60 Hz, 0.25 s, 30 deg/s, 50 and
% 100 ms, 8 s, 15 Hz, ...) are design constants of the analysis (analysis/stage2/config.py), not results.
% Sources: analysis/stage2/README.md sections 1-5, 8, 11-12; figures S2_* (core); tables describe_*, reference_*.
\section{Data collection, curation and reproducibility}\label{app:data}

This appendix documents the data decisions behind every analysis: what the app records, how the two phone
platforms are brought to one convention, which attempts become trials, what is flagged but kept, how the labels
were screened, the reference floor that every model must beat, and how the report is rebuilt. Each subsection
states the options, the evidence and the decision. Table~\ref{tab:app:data-decisions} lists the decisions at a
glance; Table~\ref{tab:notation} gives the notation.

Evaluation is always leave-one-person-out (LOSO): a model is trained on all other persons and tested on the
held-out person, one fold per person \cite{saeb2017usecase,little2017crossvalidation}. All 95\,\% confidence
intervals (CIs) of this report are person-cluster bootstrap intervals: 2000 resamples of the persons with
replacement, each resampled person bringing all of their trials \cite{saravanan2020bootstrap}. Other
abbreviations: area under the receiver operating characteristic curve (AUC), standard deviation (SD),
standard error (SE), interquartile range (IQR), root mean square (RMS).

\begin{table}[htbp]
\centering
\caption{The data decisions of this appendix at a glance. Each row is argued in the subsection named in the
last column, with its figure or table.}\label{tab:app:data-decisions}
\footnotesize
\begin{tabularx}{\linewidth}{@{}>{\raggedright\arraybackslash}p{2.5cm}>{\raggedright\arraybackslash}X
  >{\raggedright\arraybackslash}X>{\raggedright\arraybackslash}p{1.9cm}@{}}
\toprule
Decision & Options considered & Choice & Where \\
\midrule
Model inputs & every recorded stream; motion streams only & motion streams and segmentation bounds; no heading,
  magnetometer or Euler angles & \ref{app:data:app} \\
Accelerometer sign & engine; calibration; trial votes & engine sign confirmed by the trial votes & \ref{app:data:conventions} \\
Gyroscope bias & subtract the pre-GO mean; none & none & \ref{app:data:conventions} \\
Sampling grid & native timestamps; a finer grid; 60\,Hz & linear resampling to 60\,Hz & \ref{app:data:sampling} \\
Sampling gaps & interpolate all; drop every attempt with a gap; a threshold & interpolate up to 100\,ms;
  exclude a longer gap near the movement (E4d) & \ref{app:data:sampling} \\
Value precision & round iOS to Chrome's grid; native & native on both platforms & \ref{app:data:sampling} \\
Platforms & a model per platform; platform as a feature; one model & one model; platform is a reporting
  stratum & \ref{app:data:sampling} \\
Statistical unit & app code; session; person & the curated person & \ref{app:data:curation} \\
Exclusions & gates only; exclude every gate failure; exclude segments without a valid representation &
  E4b--E4d exclude; the gates are scored apart & \ref{app:data:curation}, \ref{app:data:flags} \\
Partial sessions & drop; keep & keep & \ref{app:data:completion} \\
Approach window & offline segmentation; stored online segmentation; screen contact & stored online segmentation
  (the oracle window) & \ref{app:data:segmentation} \\
Natural-use set & every free-choice trial; the first per person & the first per person (N1) & \ref{app:data:curation} \\
Suspect labels & correct them; keep them & keep them; sensitivity analyses & \ref{app:data:labels} \\
Reference floor & none; a pre-specified simple model & mirror-augmented yaw logistic regression &
  \ref{app:data:reference} \\
\bottomrule
\end{tabularx}
\end{table}

% ------------------------------------------------------------------------------------------------------------
\subsection{The app and what it records}\label{app:data:app}

The collector is a web app opened from a link or a QR code in the participant's own phone browser; nothing is
installed. On an iPhone it runs in Safari (every iOS browser uses the WebKit engine); on Android in Chrome or
Samsung Internet (the Chromium engine). The street protocol (Table~\ref{tab:protocol}) is hands-free: spoken
and on-screen prompts announce the hand, the grip and the speed, the app waits for the phone to be still, gives
a GO cue, detects the movement while it happens and signals its end. Figure~\ref{fig:gesture} shows the gesture
as the app's instructions draw it. Recruitment, compensation and a photo of a session (Figure~\ref{fig:photo}) are
in Section~\ref{sec:data}.

Every attempt is uploaded as one record as soon as it ends (Figure~\ref{fig:system}). A record holds:
\begin{itemize}[nosep]
  \item the motion streams of the browser's \code{devicemotion} event at about 60\,Hz, the browser maximum:
    the linear acceleration $\va$, the specific force $\vf$ (acceleration including gravity) and the rotation
    rate $\vomega$;
  \item orientation streams, all best effort and metadata only: Euler angles, relative and absolute
    (magnetometer-fused) quaternions, the magnetometer where a browser exposes it, and on iOS the compass
    heading;
  \item screen touches (time, position, contact size, force) and screen rotations;
  \item the app's online segmentation and the cue times: GO, onset, end, end cue
    (Appendix~\ref{app:data:segmentation});
  \item the return of the arm after the end cue, with its own onset and end (data schema 1.2);
  \item a short profile (age, sex, dominant hand, usual phone hand, optional height) and device details.
\end{itemize}
The app collects no camera, microphone or location data.

\paragraph{Decision: what feeds the models.} The options were to use every stream or the motion streams only.
Only the motion streams and the segmentation bounds reach a model. The heading streams (compass, absolute
orientation, magnetometer) are excluded for two reasons. Earth's magnetic field is not mirrored with the
movement, so a heading cannot be mirrored (Equation~\eqref{eq:mirror}), and a heading partly reveals the place
where the phone was used. The Euler angles are excluded because they hit gimbal lock at the upright start pose,
carry an arbitrary heading and, on Chrome, update only when an angle changes by at least $0.1^\circ$. Touches
serve only to describe the contact (Appendix~\ref{app:data:segmentation}). The return is never pooled with
approaches: it reverses the approach in time, and its yaw points to the other hand
(Appendix~\ref{app:data:segmentation}; Appendix~\ref{app:prod}). The profile never enters a model:
it lives in a person table, apart from the trial table, and only stratifies results.

\begin{figure}[htbp]
\centering
\begin{subfigure}[t]{0.32\linewidth}\centering
  \repoimage{gesture-a}{../docs/img/gesture/right_upright_natural_start.png}\caption{Arm extended, still}
\end{subfigure}\hfill
\begin{subfigure}[t]{0.32\linewidth}\centering
  \repoimage{gesture-b}{../docs/img/gesture/right_upright_natural_halfway.png}\caption{Mid-reach}
\end{subfigure}\hfill
\begin{subfigure}[t]{0.32\linewidth}\centering
  \repoimage{gesture-c}{../docs/img/gesture/right_upright_natural_touch.png}\caption{Nose touch and hold}
\end{subfigure}
\caption{The gesture as the app's instructions show it (right hand, upright grip, natural speed).}
\label{fig:gesture}
\end{figure}

\begin{figure}[htbp]
\centering
\begin{tikzpicture}[node distance=5mm and 6mm,>={Stealth[length=2mm]},
  box/.style={draw=black!55,rounded corners=2pt,fill=black!3,align=center,
              font=\sffamily\scriptsize,minimum height=11mm,text width=24mm}]
  \node[box] (phone) {phone\\collector web app\\uploads every attempt};
  \node[box,right=of phone] (heroku) {Heroku app (EU)\\validates records};
  \node[box,right=of heroku] (pg) {Postgres\\one row per record};
  \node[box,right=of pg] (snap) {snapshot + curation\\(persons P01--P20)};
  \node[box,below=of snap] (cache) {Stage 2 cache\\standardised trials};
  \node[box,left=of cache] (lanes) {analysis lanes\\leave-one-person-out};
  \node[box,left=of lanes] (out) {figures, tables,\\numbers, model export};
  \node[box,left=of out,dashed] (pred) {Stage 3 inference page\\(planned)};
  \draw[->] (phone) -- (heroku);
  \draw[->] (heroku) -- (pg);
  \draw[->] (pg) -- (snap);
  \draw[->] (snap) -- (cache);
  \draw[->] (cache) -- (lanes);
  \draw[->] (lanes) -- (out);
  \draw[->,dashed] (out) -- (pred);
\end{tikzpicture}
\caption{Data flow from the phone to this report. Every figure, table and number of the report is regenerated
from the saved analysis outputs by one command (Appendix~\ref{app:data:repro}).}\label{fig:system}
\end{figure}

% ------------------------------------------------------------------------------------------------------------
\subsection{Sensor conventions and platform normalisation}\label{app:data:conventions}

Table~\ref{tab:sensors} summarises how the two browser engines deliver the sensors
\cite{w3c2025deviceorientation}. Both use the device frame of the World Wide Web Consortium (W3C) specification: $x$ toward the right edge,
$y$ toward the top edge, $z$ out of the screen. Two differences matter for the analysis. WebKit inverts the
sign of every accelerometer axis, and Chrome rounds every value to 0.1 unit while WebKit does not.
Figure~\ref{fig:app:data-conventions} checks the resulting conventions end to end on real movements.

\begin{table}[htbp]
\centering
\caption{What a phone browser delivers (condensed from the project's sensor notes and the W3C specification
\cite{w3c2025deviceorientation}).}\label{tab:sensors}
\footnotesize
\begin{tabularx}{\linewidth}{@{}l>{\raggedright\arraybackslash}X>{\raggedright\arraybackslash}X@{}}
\toprule
 & iOS: WebKit (every iOS browser) & Android: Chromium (Chrome, Samsung Internet) \\
\midrule
Motion rate & $\approx$60\,Hz from a 1/60\,s timer; duplicated or skipped samples & 60\,Hz pump \\
Value precision & unrounded doubles & rounded to 0.1\,m/s$^2$ and 0.1\,deg/s \\
Accelerometer sign & all axes inverted (face up: $z \approx -9.81$) & W3C (face up: $z \approx +9.81$) \\
\code{rotationRate} & deg/s; alpha, beta, gamma about $x, y, z$; not inverted & the same \\
\code{event.interval} & $\approx$0.0167, in seconds & 16, in ms, constant \\
Permission & \code{requestPermission()} inside a tap; a Cancel sticks & Chrome 151+ also has it; granted by default \\
\bottomrule
\end{tabularx}
\end{table}

\paragraph{Decision: one versioned normalisation.} Raw records are never modified. One versioned step
(\code{norm_v2}) maps every session to the W3C convention before any feature is computed. The accelerometer
sign could come from the engine (WebKit $-1$, Chromium $+1$), from a calibration, or from the data. Street
sessions have no guided calibration, so the step takes the engine's expected sign and requires the trial votes
to confirm it. In each trial with a fixed instructed grip, the rest gravity before the onset votes for the sign
that makes it point up. Every included session's sign is supported: the street sessions by the engine sign that
the votes confirm, the pilot session by its guided calibration. Android sessions resolve to $+1$ and iOS
sessions to $-1$, as the engines predict. The \code{rotationRate} mapping (alpha, beta, gamma about $x, y, z$)
is chosen by the gravity closure defined below: the W3C mapping closes with a median of
\R{describe/rr_closure_median_deg}$^\circ$, against \R{describe/rr_runner_up_median_deg}$^\circ$ for the
runner-up mapping. Sessions too short to settle it keep the W3C default.

\paragraph{Evidence: the 64 sign patterns.} An end-to-end check uses the movements themselves. The mirror of
Equation~\eqref{eq:mirror} turns a right-hand movement into a left-hand one. In device coordinates it predicts
the sign pattern $\sigma^* = (-,+,+;\,+,-,-)$ on $(a_x, a_y, a_z;\, \omega_x, \omega_y, \omega_z)$, applied
together with the mirrored grip. For each of the $2^6 = 64$ sign patterns, Figure~\ref{fig:S2_12_sign_patterns}
measures the distance between every person's mean left-hand waveform and the sign-flipped mean right-hand
waveform. The predicted pattern $\sigma^*$ ranks \R{describe/sigma_star_rank_android} of 64 on Android and
\R{describe/sigma_star_rank_ios} on iOS. A wrong axis, a wrong sign or a failed iOS normalisation would put
another pattern first.

\paragraph{Decision: gravity frame without a gyroscope-bias correction.} The heading-free channels need the up
vector $\ghat$ during the movement. The gravity estimate (\code{gravity_v1}) takes the rest gravity before the
onset and after the end, carries each through the movement with the gyroscope and blends the two. Its closure
is the angle between the rest gravity carried to the nose and the rest gravity measured there; it tests the axis
mapping and the propagation together. The options for the gyroscope bias were to subtract the mean rate before
GO when it is below 3\,deg/s (the analysis plan: the repository's Stage 2/3 plan, docs/09, whose decisions are
cited as D1, D2, \dots) or to subtract nothing. Nothing is subtracted. Street sessions
have no still calibration window, production has no GO, and the mean rate of a hand-held phone is tremor rather
than bias. Without the correction the closure has a median of \R{describe/closure_median_deg}$^\circ$, and
\R{describe/closure_share_over_3deg} of the primary trials exceed $3^\circ$
(Figure~\ref{fig:S2_17_gravity_closure}). The median is \R{describe/android_closure_median_deg}$^\circ$ on
Android and \R{describe/ios_closure_median_deg}$^\circ$ on iOS.

\begin{figure}[htbp]
\centering
\begin{subfigure}[t]{0.49\linewidth}\centering
  \genfig{core/S2_12_sign_patterns}
  \caption{The 64 sign patterns of the device-frame channels, sorted by the distance between left-hand and
  sign-flipped right-hand mean waveforms (channels scaled to unit RMS). Circled: the conventions' mirror
  $\sigma^*$, first on both platforms.}\label{fig:S2_12_sign_patterns}
\end{subfigure}\hfill
\begin{subfigure}[t]{0.49\linewidth}\centering
  \genfig{core/S2_17_gravity_closure}
  \caption{Gyro--gravity closure of the primary trials: the angle between the rest gravity carried through the
  movement by the gyroscope and the rest gravity at the nose, with no bias correction. Line: $3^\circ$.}
  \label{fig:S2_17_gravity_closure}
\end{subfigure}
\caption{Two end-to-end checks of the sensor conventions on real movements. They support the sign
normalisation, the \code{rotationRate} mapping and the gravity frame without a gyroscope-bias correction.}
\label{fig:app:data-conventions}
\end{figure}

% ------------------------------------------------------------------------------------------------------------
\subsection{Sampling, gaps, precision and platforms}\label{app:data:sampling}

\paragraph{Decision: a uniform 60\,Hz grid.} Both platforms deliver 60\,Hz. The median sample interval is
\R{describe/android_interval_median_ms}\,ms on Android and \R{describe/ios_interval_median_ms}\,ms on iOS, and
the 99th percentile is \R{describe/android_interval_p99_ms}\,ms and \R{describe/ios_interval_p99_ms}\,ms
(Figure~\ref{fig:S2_05_sampling_gaps}, left; Table~\ref{tab:describe_sampling}). Isolated duplicated samples
are dropped (flag T11, Appendix~\ref{app:data:flags}), and the streams are resampled linearly onto a uniform
60\,Hz grid. A finer grid would add no information, and the inference page sees the same browser rate. The
low-pass filter and the analysis windows are preprocessing decisions (Appendix~\ref{app:prep}).

\paragraph{Evidence: two kinds of gap.} The intervals have a tail of gaps (Figure~\ref{fig:S2_05_sampling_gaps},
right). Android gaps longer than 50\,ms are short, with a median of \R{describe/android_gap_median_ms}\,ms. iOS
gaps are several times longer, with a median of \R{describe/ios_gap_median_ms}\,ms. In the \R{describe/ios_long_gap_attempts}
iOS attempts with a gap longer than 200\,ms, these gaps last a median of \R{describe/ios_long_gap_median_ms}\,ms and
start a median of \R{describe/ios_long_gap_start_median_s}\,s after the recording starts; only
\R{describe/ios_long_gap_share_after_cue} of them follow an audio cue within 0.4\,s. This pattern points
to a recurring stall of the collection app's main thread, not to the sensor. The inference page must keep its
main thread free while it records.

\paragraph{Decision: interpolate short gaps, exclude long gaps near the movement.} The options were to
interpolate every gap, to drop every attempt with a gap, or to set a threshold. Gaps up to 100\,ms (at most
5 samples) are interpolated linearly. A longer gap stays missing in the stream. When it falls within
$[\text{onset} - 0.5\,\text{s},\ \text{end} + 0.5\,\text{s}]$ it removes part of the movement, which
interpolation cannot restore, and the trial is excluded (rule E4d, Appendix~\ref{app:data:curation}). Elsewhere
in the recording it does no harm and stays missing. The long gaps are an iOS phenomenon, so E4d removes mostly
iOS trials (Table~\ref{tab:describe_sampling}, attempts with a gap longer than 100\,ms).

\paragraph{Decision: native precision.} Chrome's rounding to a 0.1-unit grid adds noise with an SD of
\R{describe/quantisation_noise_sd} unit ($0.1/\sqrt{12}$). The noise of a hand-held phone in the still window
before GO is far larger (Figure~\ref{fig:S2_06_noise_vs_quantisation}; medians over trials): the SD of the
specific force is \R{describe/android_rest_acc_sd}\,m/s$^2$ on Android and \R{describe/ios_rest_acc_sd}\,m/s$^2$
on iOS, and the SD of the angular velocity \R{describe/android_rest_gyro_sd}\,deg/s and
\R{describe/ios_rest_gyro_sd}\,deg/s. The options were to round iOS values to Chrome's grid, so that the
platforms match, or to keep the native values. Both platforms keep their native precision; rounding iOS is an
ablation (A3, Appendix~\ref{app:prep}).

\paragraph{Decision: one model for both platforms.} After normalisation the platforms agree where it matters
(Figure~\ref{fig:S2_14_platforms}; Table~\ref{tab:describe_platforms}). The class-mean yaw rate about gravity,
$\vomega\cdot\ghat$, separates the hands with the same sign on both. The first-half yaw (the yaw about gravity
accumulated over the first half of the approach) separates them with a pooled AUC of
\R{describe/fhy_auc_android} on Android and \R{describe/fhy_auc_ios} on iOS. iOS approaches are shorter (median
\R{describe/ios_duration_median_s}\,s against \R{describe/android_duration_median_s}\,s) and turn slightly
faster (peak $|\vomega|$ \R{describe/ios_peak_gyro_median} against \R{describe/android_peak_gyro_median}\,deg/s).
The options were a model per platform, platform as an input feature, or one model for both. One model serves
both platforms, and platform is a reporting stratum with a leave-one-platform-out check. Phones differ in their
sensors \cite{stisen2015heterogeneity}; here the normalisation removes the sign differences, and what remains
(the iOS stalls, the rounding, the durations) is reported per platform. With \R{describe/persons_ios} iOS
persons, a per-platform model would also be trained on very few people.

\begin{figure}[htbp]
\centering
\genfig{core/S2_05_sampling_gaps}
\caption{Sampling of the complete attempts. Left: the distribution of sample intervals on each platform (both
run at 60\,Hz). Right: every interval longer than 50\,ms, by platform and by where it falls relative to the
movement (``near'': within 0.5\,s of the approach). The line at 100\,ms is the threshold of rule E4d. Supports
the 60\,Hz grid, the interpolation of short gaps and the exclusion of long gaps near the movement.}
\label{fig:S2_05_sampling_gaps}
\end{figure}

\gentab{describe/describe_sampling}

\begin{figure}[htbp]
\centering
\genfig{core/S2_06_noise_vs_quantisation}
\caption{Hand-held noise in the still window before GO, per trial and platform (bar: median), against the
noise of Chrome's 0.1-unit rounding (SD \R{describe/quantisation_noise_sd}, vertical line). Left: specific
force; right: angular velocity. The rounding is far below the physiological noise floor, so native precision
is kept on both platforms.}
\label{fig:S2_06_noise_vs_quantisation}
\end{figure}

\begin{figure}[htbp]
\centering
\genfig{core/S2_14_platforms}
\caption{Android and iOS after standardisation: the class-mean yaw rate about gravity over the time-normalised
approach (band: $\pm 1$ SE across persons) and the first-half yaw per trial (bar: median). The hands separate
with the same sign on both platforms, which confirms the sign normalisation and supports one model for both.}
\label{fig:S2_14_platforms}
\end{figure}

\gentab{describe/describe_platforms}

% ------------------------------------------------------------------------------------------------------------
\subsection{Curation and the exclusion cascade}\label{app:data:curation}

The analysis reads one frozen snapshot of the upload database: \R{describe/sessions_total} sessions with
\R{describe/attempts_total} attempts, each validated against the data schema, with no duplicate and no invalid
record. The researcher's curation file maps every session to a person (P01--P20) and marks
\R{describe/sessions_included} sessions as included. The file's free-text reasons are replaced by fixed codes,
and the private file that links persons to names is never opened (Appendix~\ref{app:data:repro}).

\paragraph{Decision: the person is the statistical unit.} The options were the app's participant code, the
session, or the curated person. The analysis groups by the curated person. One app code held two persons who
shared a phone, and one person held several codes; grouping by code would put one person on both sides of a
fold. The evaluation harness refuses any grouping that splits a person. \R{describe/n_persons_multi_session}
persons contributed several sessions; one of these sessions is the pilot of the earlier protocol (1.0.0). The
demographics are read from each session's own record, agree across each person's sessions, and are stored in
the person table only, so they cannot reach a model.

\paragraph{The exclusion cascade.} Rules apply in order, and an attempt is charged to the first rule it fails
(Table~\ref{tab:describe_exclusion_flow}). Rule E3 diverts rather than removes.
\begin{itemize}[nosep]
  \item \textbf{E0} removes the \R{describe/excl_e0_removed} attempts of the sessions the researcher excluded:
    an incomplete run superseded by the person's later run, and runs that did not work.
  \item \textbf{E0s} (simulated records) and \textbf{E2} (the screen-away grip, collected only by the pilot
    protocol) remove nothing. Street visits never collect screen-away trials; the analysis synthesises them by
    a screen flip of the arrays (Appendix~\ref{app:prep}).
  \item \textbf{E1} removes the \R{describe/excl_e1_removed} attempts the app did not accept
    (Table~\ref{tab:describe_e1_reasons}). Most showed no movement within the time limit after GO, or a phone
    still moving at the limit; the few others were redone, rejected for a wrong grip, aborted, flagged or
    interrupted. Overall, \R{describe/attempts_per_accepted} attempts were needed per accepted trial.
  \item \textbf{E3} moves the \R{describe/free_choice_trials} free-choice trials to the natural-use set. Their
    hand was chosen and reported by the person, not instructed, so they are never trained on and are scored
    apart.
  \item \textbf{E4a} (no approach window) removes nothing.
  \item \textbf{E4b} removes \R{describe/excl_e4b_removed} approaches shorter than 0.25\,s: a reach from arm's
    length cannot be that short, so a fragment was segmented (Figure~\ref{fig:S2_08_durations_peaks}).
  \item \textbf{E4c} removes \R{describe/excl_e4c_removed} attempts whose phone did not turn (peak $|\vomega|$
    below 30\,deg/s, the app's own onset floor): every reach rotates the forearm. The natural-use set also
    loses a trial here.
  \item \textbf{E4d} removes \R{describe/excl_e4d_removed} attempts with a gap longer than 100\,ms near the
    movement (Appendix~\ref{app:data:sampling}).
  \item \textbf{E4e} (page hidden, rate below 50\,Hz, no gyroscope, unsupported gravity sign) removes nothing.
  \item \textbf{N1} keeps one natural-use trial per person: the earliest free-choice trial that passed every
    rule, the person's first unpractised movement with a hand of their choice. The
    \R{describe/natural_use_repeat_n} free-choice trials of repeat visits are reported apart.
\end{itemize}

\paragraph{The analysis sets.} The primary set holds \R{describe/primary_trials} trials of
\R{describe/n_persons} persons in \R{describe/sessions_included} sessions: \R{describe/primary_left} left-hand
and \R{describe/primary_right} right-hand trials, \R{describe/primary_android} on Android and
\R{describe/primary_ios} on iOS. Each person contributes both hands, with \R{describe/primary_per_person_min} to
\R{describe/primary_per_person_max} trials (median \R{describe/primary_per_person_median}). The natural-use set
holds \R{describe/natural_use_n} trials, one per person, and all are right-hand trials
(\R{describe/natural_use_right} right, \R{describe/natural_use_left} left). Natural-use accuracy is therefore
right-hand recall on the free-choice trials (all right-hand so far; a constant right answer would score
100\,\%). The set supports an accuracy with a CI and nothing that needs both hands, such as an AUC or a
calibration.

\paragraph{Decision: exclude only what has no valid representation.} The analysis plan (docs/09, D5) foresaw no
new exclusions: movement checks would act as predictor gates at evaluation. The options were to keep that plan,
to exclude every gate failure, or to exclude only segments without a valid representation. Rules E4b, E4c and
E4d exclude, because a fragment, a phone that does not turn and a hole in the movement have no valid
representation. The remaining gate failures are valid movements and stay in training and in the headline
(Appendix~\ref{app:data:flags}). The cache keeps the arrays of every excluded complete attempt, so a
sensitivity analysis that re-admits them needs no rebuild.

\gentab{describe/describe_exclusion_flow}

\gentab{describe/describe_e1_reasons}

% ------------------------------------------------------------------------------------------------------------
\subsection{Completion, conditions and durations}\label{app:data:completion}

\paragraph{Decision: keep partial sessions.} Of the \R{describe/street_visits} street visits,
\R{describe/visits_completed} were completed and \R{describe/visits_ended_early} ended early
(Figure~\ref{fig:S2_03_completion_durations}, left). The optional bonus tier was taken in
\R{describe/bonus_taken_completed} of the completed visits. The options were to drop partial sessions or to
keep them. They are kept: the plan balances the hands within every tier, so even a visit that stops early
contributes both hands. Dropping partial sessions would lose trials, and a whole person whose only visit
ended early.

\paragraph{Evidence: visits run long.} A completed visit took a median of \R{describe/visit_minutes_median}\,min
(IQR \Rb{describe/visit_minutes_q1}--\Rb{describe/visit_minutes_q3}\,min), not the planned 5\,min
(Figure~\ref{fig:S2_03_completion_durations}, right). The tiers took medians of
\R{describe/tier1_minutes_median}, \R{describe/tier2_minutes_median} and \R{describe/tier3_minutes_median}\,min
(tiers 1 to 3) and \R{describe/tier4_minutes_median}\,min for the bonus tier. The next collection should plan
for the measured duration, not the app's estimate.

\paragraph{Evidence: the conditions are balanced between hands.} The primary trials span
\R{describe/n_conditions} instructed conditions of grip, speed and movement variant
(Figure~\ref{fig:S2_04_conditions}). Every condition has trials of both hands in similar numbers, so a
condition cannot stand in for the label. Upright standard reaches at natural speed dominate, and the rare
variants hold too few trials for a per-condition accuracy; the secondary analyses pool them
(Appendix~\ref{app:sec}).

\paragraph{Evidence: durations and the thresholds of E4b and E4c.} The oracle approach (defined in
Appendix~\ref{app:data:segmentation}) lasts a median of \R{describe/duration_median_s}\,s, and
\R{describe/duration_share_le_2p5s} of approaches last at most 2.5\,s (Figure~\ref{fig:S2_08_durations_peaks};
Table~\ref{tab:describe_durations}). The instructed speed separates durations only loosely: the medians are
\R{describe/duration_fast_s}\,s for fast, \R{describe/duration_natural_s}\,s for natural and
\R{describe/duration_slow_s}\,s for slow movements. Duration is therefore a poor proxy for the instructed speed. The reaction time from GO to the onset has a median of \R{describe/reaction_median_ms}\,ms.
The fixed analysis window ($-0.5$ to $+2.5$\,s around the onset) and the time normalisation follow from these
durations (Appendix~\ref{app:prep}).

\begin{figure}[htbp]
\centering
\genfig{core/S2_03_completion_durations}
\caption{Completion. Left: accepted attempts per tier and street session (pseudonymous session labels; \R{describe/visits_ended_early}
visits ended early). Right: wall-clock durations of the setup, of each tier and of the whole visit, for the
completed visits (bar: median). Partial sessions still hold both hands in every tier, so they are kept.}
\label{fig:S2_03_completion_durations}
\end{figure}

\begin{figure}[htbp]
\centering
\genfig{core/S2_04_conditions}
\caption{Primary trials per instructed condition (grip $\cdot$ speed $\cdot$ movement variant $\cdot$ start
position), left against right hand, on a log scale. Every condition has both hands in similar numbers, so
no condition can stand in for the label.}
\label{fig:S2_04_conditions}
\end{figure}

\begin{figure}[htbp]
\centering
\genfig{core/S2_08_durations_peaks}
\caption{Left: oracle approach duration by instructed speed and platform (bar: median). Solid line: the
0.25\,s threshold of E4b (a lower bound of gate G2); dashed line: the 8\,s upper bound of gate G2 (longer
approaches stay in the primary set and are scored apart). Right: peak angular speed within the approach, with
the 30\,deg/s threshold of E4c. Supports excluding fragments and phones that do not turn.}
\label{fig:S2_08_durations_peaks}
\end{figure}

\gentab{describe/describe_durations}

% ------------------------------------------------------------------------------------------------------------
\subsection{Segmentation: the oracle window, online and offline}\label{app:data:segmentation}

Two segmentations exist for every recording. The \emph{online} segmentation is computed while the movement
happens, from the samples seen so far: the app's \code{online_v1} detector drives the end cue and stores its
bounds with the attempt. Its onset is back-dated to the end of the stillness before the movement, and its end
to the start of the final stillness at the nose; the end cue follows 250\,ms after the end. The
\emph{offline} segmentation is computed after the recording, on the whole stream: \code{offline_v1} re-segments
each recording in the analysis and needs some stillness after the end to find it. The \emph{oracle window} is
the approach between the stored online bounds. In production nothing marks the movement, and a cue-free
detector must find it (Appendix~\ref{app:prod}).

\paragraph{Decision: the oracle window is the stored online segmentation.} The options were the offline
segmentation (the analysis plan, docs/09 D10), the stored online segmentation, or the first screen contact as
the end of the approach. The evidence is in Figure~\ref{fig:S2_09_segmentation_contact}. The offline
segmentation finds no end in \R{describe/offline_no_end_share_street} of street attempts, because the arm
returns soon after the end cue and leaves no long stillness at the nose. Where both exist they agree closely at
the onset and less closely at the end; \R{describe/flag_seg_disagree_100ms} primary trials have onsets more than
100\,ms apart (flag T9). The stored online bounds exist for every complete attempt and can be reproduced
exactly from the record. Screen contact is unreliable as a bound: only \R{describe/touch_share_toward} of the
screen-toward attempts record a touch, and the first touch falls at a median \R{describe/touch_position_median}
of the oracle window, \R{describe/touch_before_end_ms}\,ms before its end; the last part of the window is the
phone settling against the nose. The offline bounds stay in the cache for an ablation (A5); the sensitivity of the models to the window
bounds is tested by jittering them (Appendix~\ref{app:prep}).

\paragraph{Evidence: what a trial looks like.} Figure~\ref{fig:S2_10_example_traces} shows one left-hand and
one right-hand trial of the same person, phone, grip and speed, after standardisation. The yaw rate about
gravity takes opposite signs early in the two approaches. In both trials the return of the arm starts less than a second
after the end cue and is as large as the approach. A detector that hands the return to an approach model
gets a reversed movement and a confidently wrong answer; the return is therefore labelled separately and never
pooled with approaches (Appendix~\ref{app:prod}). A clean return window exists for \R{describe/return_trials}
of the primary trials.

\begin{figure}[htbp]
\centering
\genfig{core/S2_09_segmentation_contact}
\caption{Left: offline minus stored online segmentation bounds, at the onset and at the end (clipped at
$\pm$800\,ms; the offline segmentation finds no end in \R{describe/offline_no_end_share_street} of street
attempts). Right: the first screen touch as a fraction of the oracle window (0: onset, 1: end) in the
screen-toward attempts. Supports the stored online bounds as the oracle window, with contact not used as a
bound.}
\label{fig:S2_09_segmentation_contact}
\end{figure}

\begin{figure}[htbp]
\centering
\genfig{core/S2_10_example_traces}
\caption{One left-hand and one right-hand trial of P15 (same phone, upright grip, natural speed),
standardised: device-frame angular velocity and linear acceleration (top rows) and the heading-free
gravity-frame yaw rate and tilt rate (bottom). Grey band: the oracle approach; beige band: the return of the
arm. The yaw rate has opposite signs early in the two approaches; the return is a large separate movement that
must be kept apart.}
\label{fig:S2_10_example_traces}
\end{figure}

% ------------------------------------------------------------------------------------------------------------
\subsection{Flags and predictor gates}\label{app:data:flags}

Two kinds of trial-level column are computed for every trial and never remove one
(Figure~\ref{fig:S2_18_flags_gates}; Table~\ref{tab:describe_flags_gates}).

\paragraph{Flags} (the trial checks T8--T12 of the data schema, docs/02, and a screen-rotation warning) describe
the recording and serve only to stratify results:
\begin{itemize}[nosep]
  \item T8, grip not as instructed (\R{describe/flag_grip_mismatch} trials): the grip before the onset differs
    from the instructed one, by the app's check or by the standardised up vector. The hand label is kept,
    because the hand instruction concerns the hand, not the grip.
  \item T9, online and offline onsets more than 100\,ms apart (\R{describe/flag_seg_disagree_100ms} trials).
  \item T10, possible clipping (\R{describe/flag_clip_candidate} trials): $|\vf|$ above 3.5\,g or $|\vomega|$
    above 950\,deg/s somewhere in the raw recording.
  \item T11, repeated samples (\R{describe/flag_repeated_samples} trials): isolated duplicated samples were
    dropped before resampling.
  \item T12, a false start before GO (\R{describe/flag_false_start} trials): the hold restarted before GO.
  \item Screen rotated during the movement (\R{describe/flag_screen_rotated} trials).
\end{itemize}

\paragraph{Predictor gates} are the checks the inference page will apply before it answers (docs/09 D5). Gate
G1 requires a rate of at least 50\,Hz, no gap longer than 50\,ms within 0.3\,s of the approach, a visible page
and a gyroscope; \R{describe/gate_g1_failing} primary trials fail it, all on Android, with gaps between 50 and
100\,ms that were interpolated. Gate G2 requires a duration between 0.25 and 8\,s and a peak $|\vomega|$ of at
least 30\,deg/s; \R{describe/gate_g2_failing} trials fail it, nearly all instructed slow reaches longer than
8\,s. In all, \R{describe/gate_pass_failing} primary trials fail G1 or G2.

\paragraph{Decision: gate failures are kept and scored apart.} The plan would not score gate failures. The
options were to drop them from training and evaluation, or to keep them and report the gated metrics as well.
They are kept in training and in the headline metrics, and every result also reports its metrics on the trials
that pass both gates, as the inference page would see them. The headline is then the harder, conservative
number. For the reference floor (Appendix~\ref{app:data:reference}) the accuracy is
\R{reference/accuracy} on all primary trials and \R{reference/gated_accuracy} on the trials that pass both
gates.

\begin{figure}[htbp]
\centering
\genfig{core/S2_18_flags_gates}
\caption{Flags (T8--T12, screen rotation) and predictor gates (G1, G2) in the primary trials, by platform.
Flags are reported and used for stratification; trials failing a gate stay in training and in the headline,
and every result also reports the metrics on the trials that pass both gates.}
\label{fig:S2_18_flags_gates}
\end{figure}

\gentab{describe/describe_flags_gates}

% ------------------------------------------------------------------------------------------------------------
\subsection{Label consistency and the swap signature}\label{app:data:labels}

Hands-free collection has a blind spot. An accepted attempt only means ``completed and not redone'', so a
person who used the wrong hand for a whole block leaves no trace in the record. The screen in
Figure~\ref{fig:S2_20_label_consistency} checks each of the \R{describe/label_blocks} planned blocks of the
primary trials: the share of its trials whose first-half yaw sign agrees with the instructed hand.

\paragraph{The swap signature.} A hand swap at a block change produces a specific pattern: two adjacent blocks of
opposite instructed hands, each of at least 3 trials, in which every trial contradicts its label. A person
whose hand moves differently produces another pattern, in which one hand disagrees and the other agrees. The
signature is found from the data, without a typed person or block. It appears in \R{describe/swap_signature_blocks}
blocks: blocks 1 and 2 of P08, in which no trial agrees with its label. Low agreement confined to one hand
of a person is not the signature. A post hoc example is P17, a left-hander, whose left-hand blocks mostly
disagree with the rule while the right-hand blocks agree. Handedness is examined in Appendix~\ref{app:sec}.

\paragraph{Decision: the labels are never changed.} The options were to correct suspected labels or to keep the
instructed labels and test their influence. The labels stay as instructed, and the ringed blocks are flagged for
review. Two safeguards follow. A sensitivity analysis drops the ringed blocks from training, or
from training and test (Appendix~\ref{app:data:reference}). The wrong-hand screen of the analysis plan
(docs/09 D6) runs inside each training fold: an antisymmetric logistic model on the chiral features, scored by
inner LOSO over the training persons, flags training trials with $P(\text{label}) < 0.05$ and runs of at least 3
consecutive trials with $P(\text{label}) < 0.25$. Flagged trials leave training only, never evaluation, so the
cleaning cannot inflate accuracy. Applied once to the whole primary set, with LOSO scores, these rules flag
\R{describe/screen_flagged_trials} trials.

\begin{figure}[htbp]
\centering
\genfig{core/S2_20_label_consistency}
\caption{Label consistency of every planned block of every session (dot size: trials; colour: instructed hand;
fill: share of trials whose first-half yaw sign agrees with the instructed hand; text: agreeing/total). Ringed:
the swap signature, two adjacent blocks of opposite hands in which every trial contradicts its label. The
labels are kept; the ringed blocks enter a sensitivity analysis.}
\label{fig:S2_20_label_consistency}
\end{figure}

% ------------------------------------------------------------------------------------------------------------
\subsection{The reference floor and label-noise sensitivity}\label{app:data:reference}

\paragraph{Decision: a pre-specified reference floor.} Every model is compared against one simple,
pre-specified model, defined with the analysis foundation before any modelling lane ran. It is a logistic regression on the cumulative yaw about gravity
at 8 instants of the time-normalised approach, trained with mirror augmentation (each training trial is added
again mirrored, with the opposite label). It is a sanity floor, not a model choice; lane m4 uses the same model,
under the name yaw8, as the hand classifier of its production simulation (Appendix~\ref{app:prod}). Its LOSO accuracy is
\Rci{reference/accuracy}, with an AUC of \R{reference/roc_auc} and a log loss of \R{reference/log_loss}
(Table~\ref{tab:reference_models}). On the natural-use set it scores \R{reference/natural_use_accuracy}: this is
right-hand recall on the free-choice trials (all right-hand so far; a constant right answer would score
100\,\%).

\paragraph{Evidence: several attempts help.} When the \pright{} of $k$ attempts of one person and hand are
averaged, the per-person balanced accuracy rises from \R{reference/acc_k1} at $k = 1$ to \R{reference/acc_k3} at
$k = 3$ and \R{reference/acc_k5} at $k = 5$ (Figure~\ref{fig:S2_21_reference_sensitivity}, left;
Table~\ref{tab:reference_k_attempts}). Every $k$ uses the same \R{reference/k_n_persons} persons, those with at
least 5 trials of each hand. The inference page should therefore ask for several movements.

\paragraph{Evidence: label noise does not move the floor.} Three variants test the labels
(Figure~\ref{fig:S2_21_reference_sensitivity}, right; Table~\ref{tab:reference_models}). Each is paired against
the reference on the same trials, with an exact sign-flip test over persons; the family is these three
comparisons. Without the swap-signature blocks in training the accuracy changes by
\Rb{reference/swap_blocks_train_only_diff}~points (unadjusted $p = \R{reference/swap_blocks_train_only_p}$).
Without them in training and test it changes by \Rb{reference/swap_blocks_removed_diff}~points
($p = \R{reference/swap_blocks_removed_p}$). With the D6 screen in every fold it changes by
\Rb{reference/d6_screen_diff}~points ($p = \R{reference/d6_screen_p}$). None differs even before the Holm
adjustment over the three, which can only raise these $p$-values. The labels therefore stay as instructed, and
the D6 screen stays an optional sensitivity analysis.

\begin{figure}[htbp]
\centering
\genfig{core/S2_21_reference_sensitivity}
\caption{Reference floor. Left: per-person balanced accuracy when the \pright{} of $k$ attempts of one person
and hand are averaged (the same \R{reference/k_n_persons} persons at every $k$). Right: the label-noise
variants, paired against the reference on the same trials (95\,\% CI of the paired difference). Supports asking
for several attempts and keeping the instructed labels.}
\label{fig:S2_21_reference_sensitivity}
\end{figure}

\gentab{reference/reference_models}

\gentab{reference/reference_k_attempts}

% ------------------------------------------------------------------------------------------------------------
\subsection{Demographics}\label{app:data:demographics}

The \R{describe/n_persons} persons are young adults (Figure~\ref{fig:S2_01_demographics}): age
\R{describe/age_mean}\,$\pm$\,\R{describe/age_sd} years (range \R{describe/age_min} to \R{describe/age_max}),
\R{describe/n_female} women and \R{describe/n_male} men, and height \R{describe/height_mean}\,$\pm$\,\R{describe/height_sd}\,cm
(\R{describe/height_n} persons gave it). \R{describe/n_right_handed} are right-handed and
\R{describe/n_left_handed} left-handed (P11 and P17, self-report). The usual phone hand is the right for
\R{describe/n_phone_hand_right} persons and both for \R{describe/n_phone_hand_both}. \R{describe/persons_android}
persons used Android phones (\R{describe/sessions_android} sessions, \R{describe/n_device_models} device models
across all included sessions) and \R{describe/persons_ios} used iPhones (\R{describe/sessions_ios} sessions).

\paragraph{Consequences.} With \R{describe/n_left_handed} left-handers, every handedness effect is exploratory, and the left-handers
are shown individually (Appendix~\ref{app:sec}). The narrow age range limits any claim about older users.
Demographics never enter a model; they only stratify results, and the demographics add-on test is in
Appendix~\ref{app:sec}.

\begin{figure}[htbp]
\centering
\genfig{core/S2_01_demographics}
\caption{The \R{describe/n_persons} persons of the analysis: age, height by sex (\R{describe/height_n} persons
gave a height), dominant hand and usual phone hand (self-report). With \R{describe/n_left_handed} left-handers,
handedness effects can only be exploratory.}
\label{fig:S2_01_demographics}
\end{figure}

% ------------------------------------------------------------------------------------------------------------
\subsection{Reproducibility, provenance and privacy}\label{app:data:repro}

\paragraph{One command rebuilds everything.} One orchestrator regenerates every result, figure, table and
number of this report from the frozen snapshot (\code{analysis/stage2/pipeline.py}):
\begin{small}
\begin{verbatim}
cd analysis      # python = the frozen environment, .venv/Scripts/python.exe
python -m stage2.pipeline all --force   # cache, every unit, numbers, report
python -m stage2.pipeline all           # only what is missing or stale
python -m stage2.pipeline figures       # redraw every figure (no fitting)
python -m stage2.pipeline tables        # booktabs tables and CSV files
python -m stage2.pipeline numbers       # numbers files -> numbers.tex
python -m stage2.pipeline report        # asset step, then latexmk
\end{verbatim}
\end{small}
Each target takes the options \texttt{-{}-lanes} (units) and \texttt{-{}-only} (names). The Python environment is
frozen in a requirements file (\code{requirements-stage2.txt}), and every random choice draws from one fixed
seed.

\paragraph{Compute is separate from presentation.} A compute step runs the experiments and saves everything the
presentation needs: predictions, per-fold information, curves and the data behind each figure. The figures,
tables and numbers steps only read those files and never fit a model, so redrawing a figure after a change of
an axis takes seconds. To change a figure, edit the presentation defaults of its figures module, run
\texttt{python -m stage2.pipeline figures -{}-lanes <unit> -{}-only <NAME>} and rebuild the report; the asset step
copies the new file.

\paragraph{Every number is generated.} Each unit writes its numbers to a file of the form
\texttt{\{key: \{value, format, source, description\}\}}, where \code{source} names the saved file the value was
read from. The asset step turns them into the \LaTeX{} macros that print every result in this report; no result
is typed by hand. A missing key prints in red, and a final build stops on it.

\paragraph{Staleness and provenance.} The cache records its data identity: a hash over the snapshot, the
curation file, the cache parameters and the code that builds the cache. Every result records the data identity
of the cache it was made from, with the git commit and a fingerprint of the analysis code. A changed snapshot or
curation makes every result stale, and a plain \code{all} recomputes them; an edit of presentation code makes
nothing stale. The footer of every page prints the provenance of the build: \provenanceline.

\paragraph{Adding a participant.} Five steps add a person:
\begin{enumerate}[nosep]
  \item put the new session in the snapshot, or export a new snapshot and point the configuration at it;
  \item add its row to the curation file with the next person id (e.g.\ P21) and \code{include: true};
  \item run \code{python -m stage2.pipeline all}: the data identity has changed, so the cache is rebuilt and
    every result recomputed;
  \item check that \texttt{python -m stage2.scripts.expected\_counts -{}-without P21} reproduces the pinned counts
    exactly, so every change is explained by the new person, then pin the new counts with \texttt{-{}-write};
  \item run the test suite (\code{python -m pytest tests stage2/tests}).
\end{enumerate}
The last person added, P20, passed this check: without P20 every count of the previous snapshot is reproduced
exactly.

\paragraph{Privacy.} The repository is public, so no name reaches any output. Persons appear only as P01--P20,
and sessions only as pseudonymous labels such as \code{P14-s3}. The private file that links persons to names is
never opened, and the curation's free-text reasons are replaced by fixed codes. Report-facing tables refuse app
codes, backend ids, timestamps, operating-system and browser versions and screen geometry. The tests scan every
table, JSON file, \LaTeX{} file, caption, figure text, PNG and PDF, the cache tables and every source file for
the name tokens of the raw records, case-insensitively. Exploratory participants such as PD01 live outside the
snapshot and the curation; only the case-study lane reads them, no model is trained, tuned or scored on them,
and a guard test checks that none of their ids reaches the main data (Appendix~\ref{app:pd01}).

\begin{table}[htbp]
\centering
\caption{Notation.}\label{tab:notation}
\footnotesize
\begin{tabularx}{\linewidth}{@{}l>{\raggedright\arraybackslash}X@{}}
\toprule
Symbol & Meaning \\
\midrule
$\vf$, $\va$ & specific force (acceleration including gravity) and linear acceleration, device frame, m/s$^2$ \\
$\vomega$ & angular velocity (\code{rotationRate} as $x, y, z$), deg/s \\
$\ghat$ & unit vector pointing up (gravity frame), estimated by gyroscope propagation from the rest pose \\
$\Smat$, $\mathbf{M}$ & $\operatorname{diag}(-1, 1, 1)$ on the device and on the body side \\
$\Qmat$ & a grip: a constant proper rotation of the device frame \\
$m(\cdot)$ & the left--right mirror of a trial or of a feature vector \\
$\vomega\cdot\ghat$ & yaw rate about gravity (chiral, or mirror-odd: it changes sign under the mirror) \\
$\pright$ & the model's probability that the right hand made the movement \\
$\tauk$ & the scale of the summed logits $z_i$ of $k$ attempts, $\pright = \sigma(\sum_i z_i / \tauk)$, $\sigma$ the logistic sigmoid \\
\bottomrule
\end{tabularx}
\end{table}
